\begin{proposition}\label{prop:2.2.5}
设 $q:Y\to X$ 是 $X$ 的覆盖。对每个 $y\in q^{-1}(\{x\})$，存在唯一的映射
$t_y:(\tilde X,\tilde x)\to(Y,y)$，使得 $q\circ t_y=p$。给定
$f\in\Aut(\tilde X/X)^{\mathrm{op}}$，令
\begin{equation}\label{eq:2.10}
f\curvearrowright y:=(t_y\circ f)(\tilde x).
\end{equation}
这定义了 $\Aut(\tilde X/X)^{\mathrm{op}}$ 在 $q^{-1}(\{x\})$ 上的左作用。此外，若
$\phi:Y\to Z$ 是 $X$ 的两个覆盖 $q:Y\to X$ 与 $r:Z\to X$ 之间的同态，则诱导映射
$\phi_x:q^{-1}(\{x\})\to r^{-1}(\{x\})$ 是 $\Aut(\tilde X/X)^{\mathrm{op}}$-等变的。
最后，若 $\gamma\in\pi_1(X;x)$ 在命题 2.2.4 的同构下满足 $\Psi(\gamma)=f$，则
$f\curvearrowright y=\gamma\curvearrowright y$，其中 $\gamma\curvearrowright y$ 表示
(2.5) 中定义的单值作用。
\end{proposition}

\begin{proof}
先注意 $f\curvearrowright y$ 确实属于 $q^{-1}(\{x\})$。计算得
\begin{equation}\label{eq:2.11}
q(f\curvearrowright y)
=q\bigl((t_y\circ f)(\tilde x)\bigr)
=(q\circ t_y)\circ f(\tilde x)
=(p\circ f)(\tilde x)=p(\tilde x)=x.
\end{equation}
其次有 $t_y\circ f=t_{f\curvearrowright y}$。事实上，
$q\circ(t_y\circ f)=(q\circ t_y)\circ f=p\circ f=p$，且按 $f\curvearrowright y$ 的定义
$(t_y\circ f)(\tilde x)=f\curvearrowright y$，故由 $t_{f\curvearrowright y}$ 的唯一性得
$t_y\circ f=t_{f\curvearrowright y}$。特别地，
\[
(f\cdot g)\curvearrowright y=(g\circ f)\curvearrowright y
=\bigl(t_y\circ(g\circ f)\bigr)(\tilde x)
=\bigl((t_y\circ g)\circ f\bigr)(\tilde x)
=\bigl(t_{g\curvearrowright y}\circ f\bigr)(\tilde x)
=f\curvearrowright(g\curvearrowright y),
\]
故确实有 $\Aut(\tilde X/X)^{\mathrm{op}}$ 在 $q^{-1}(\{x\})$ 上的左作用。映射
$\phi_x:q^{-1}(\{x\})\to r^{-1}(\{x\})$ 的等变性来自下图的可交换性：
\[
\begin{tikzcd}[column sep=large, row sep=large]
 & \tilde X \arrow[dl, "t_y"'] \arrow[dr, "t_{\phi(y)}"] & \\
 Y \arrow[dr, "q"'] \arrow[rr, "\phi"] & & Z \arrow[dl, "r"]\\
 & X &
\end{tikzcd}
\]
我们要证明 $\phi\circ t_y=t_{\phi(y)}$（注意 $r\circ t_{\phi(y)}=q\circ t_y$，因为二者都
等于 $p:X̃\to X$；又 $r\circ\phi=q$，因为按假定 $\phi$ 是覆盖的同态）。事实上，
$(\phi\circ t_y)(\tilde x)=\phi(t_y(\tilde x))=\phi(y)=t_{\phi(y)}(\tilde x)$，且
$r\circ(\phi\circ t_y)=(r\circ\phi)\circ t_y=q\circ t_y=p$，故由 $t_{\phi(y)}$ 的唯一性
\begin{equation}\label{eq:2.12}
\phi\circ t_y=t_{\phi(y)}.
\end{equation}
于是
\[
\phi(f\curvearrowright y)
=\phi\bigl((t_y\circ f)(\tilde x)\bigr)
=(\phi\circ t_y)\bigl(f(\tilde x)\bigr)
=(t_{\phi(y)}\circ f)(\tilde x)
=f\curvearrowright\phi(y),
\]
故 $\phi_x:=\phi|_{q^{-1}(\{x\})}$ 确实是 $\Aut(\tilde X/X)^{\mathrm{op}}$-等变的。最后，
由命题 2.2.4，若 $f=\Psi(\gamma)$ 则 $f(\tilde x)=\gamma\curvearrowright\tilde x$。由于
$t_y$ 对单值作用是 $\pi_1(X;x)$-等变的（见 (2.7)），故
\[
f\curvearrowright y=(t_y\circ f)(\tilde x)
=t_y(\gamma\curvearrowright\tilde x)
=\gamma\curvearrowright t_y(\tilde x)
=\gamma\curvearrowright y.
\]
\end{proof}

$\Aut(\tilde X/X)^{\mathrm{op}}$ 在万有覆盖 $p:\tilde X\to X$ 的 $x$ 处纤维上的左作用，
与 $\Aut(\tilde X/X)^{\mathrm{op}}$ 在 $\tilde X$ 上的右作用之间的关系，见习题 2.2.27。

在下文中，我们将继续把 $X$ 的单连通万有覆盖的存在性当作已知，并主要把
$\pi_1(X;x)$ 想成群 $\Aut(\tilde X/X)^{\mathrm{op}}$；要点在于它右作用于 $\tilde X$，而
左作用于 $X$ 的任一覆盖在 $x$ 处的纤维。对我们来说，记号 $\pi_1(X;x)$ 比
$\Aut(\tilde X/X)^{\mathrm{op}}$ 更好，因为它反映了 $X$ 中基点的选取，这一点在考虑实
曲线时将特别重要（见命题 2.2.13 与命题 2.2.17）。从现在起，我们用 $f$ 表示
$\pi_1(X;x)$ 中的元素，并把 $f$ 对 $y\in\tilde X$ 的作用记作 $y\curvearrowleft f$。
有了这些准备，我们就能从任一左 $\pi_1(X;x)$-集 $F$ 出发构造 $X$ 的一个覆盖，只需令
\[
Y:=\tilde X\times_{\pi_1(X;x)}F:=\pi_1(X;x)\backslash(\tilde X\times F),
\]
其中商空间是对 $\tilde X\times F$ 上的左 $\pi_1(X;x)$-作用取的，该作用对任意
$f\in\pi_1(X;x)$ 与 $(y,v)\in\tilde X\times F$ 规定为
\[
f\curvearrowright(y,v)=(y\curvearrowleft f^{-1},\,f\curvearrowright v).
\]
这一构造给出纤维函子（以 $\pi_1(X;x)$ 为其自同构群）的一个准逆，从而得到下述对拓扑
覆盖空间的 Galois 理论式处理，Grothendieck 在 [6] 中首次如此表述。

\begin{theorem}[拓扑覆盖空间的 Galois 理论]\label{thm:2.2.6}
设 $X$ 是连通、Hausdorff 且半局部单连通的拓扑空间。$X$ 的 Hausdorff 拓扑覆盖空间
范畴 $\mathrm{Cov}_X$ 等价于离散左 $\pi_1(X;x)$-集范畴。一对显式的准逆函子由「$x$ 处
纤维」函子
\[
\mathrm{Fib}_x:\ \mathrm{Cov}_X\longrightarrow\pi_1(X;x)\text{-}\mathbf{Sets},
\qquad (q:Y\to X)\longmapsto q^{-1}(\{x\})
\]
与函子
\[
Q:\ \pi_1(X;x)\text{-}\mathbf{Sets}\longrightarrow\mathrm{Cov}_X,
\qquad F\longmapsto\tilde X\times_{\pi_1(X;x)}F
\]
给出。映射
\[
\{\pi_1(X;x)\text{ 的子群}\}\longrightarrow\{\text{传递 }\pi_1(X;x)\text{-集}\},
\qquad H\longmapsto F:=\pi_1(X;x)/H
\]
给出 $\pi_1(X;x)$ 的子群与 $X$ 的连通覆盖之间的一个双射。在此双射下，$X$ 的 Galois
覆盖（即连通覆盖 $q:Y\to X$，使得 $\Aut(Y/X)$ 在 $q$ 的任一给定纤维上的作用传递）对应
于 $\pi_1(X;x)$ 的正规子群。
\end{theorem}

特别地，覆盖 $q:Y\to X$ 连通，当且仅当 $q^{-1}(\{x\})$ 是传递的
$\pi_1(X;x)$-集；它是 Galois 的，当且仅当 $q^{-1}(\{x\})$ 是传递的
$\pi_1(X;x)$-集，并且带有相容的群结构（即带有使该集合成为 $\pi_1(X;x)$ 的商群的群
结构）。定理 2.2.6 的证明例如可参看 [18]。注意重要的是知道函子 $\mathrm{Fib}_x$ 建立了
$X$ 的覆盖与左 $\pi_1(X;x)$-集之间的范畴等价，而万有覆盖其实只被用来构造
$\mathrm{Fib}_x$ 的一个准逆。事实上，这正是可以推广到实曲线覆盖的那一部分，我们在
2.2.2 小节将看到这一点。另一个有趣之处是：若只考虑 $\mathrm{Fib}_x$ 在 $X$ 的连通
Hausdorff 覆盖所成满子范畴上的限制，则由 $\tilde X$ 的泛性质有
$\mathrm{Fib}_x(Y)=\mathrm{Hom}_{\mathrm{Cov}_X}(\tilde X;Y)$。换言之，函子
$\mathrm{Fib}_x$ 被万有覆盖 $(\tilde X,\tilde x)$ 表出（见习题 2.2.26）。

最后，若 $X$ 是黎曼面、$p:Y\to X$ 是拓扑覆盖映射且 $Y$ 是 Hausdorff 空间，则 $Y$ 上
存在唯一的全纯结构使 $p$ 成为全纯映射（例如可参看 [5, 18]）。因此，连通黎曼面 $X$ 的
Hausdorff 拓扑覆盖范畴等价于 $X$ 的全纯覆盖空间范畴。特别地，$X$ 在黎曼面范畴中有
万有覆盖，且拓扑基本群 $\pi_1(X;x)=\Aut(\tilde X/X)^{\mathrm{op}}$ 通过全纯变换作用在
$\tilde X$ 上，并在定理 2.2.6 的意义下分类 $X$ 的全纯覆盖。我们以后会看到，这一观点对
克莱因曲面有确切的类比；见 2.2.3 小节。

\subsection{克莱因曲面的基本群}
\label{sec:02.2.2}

由定理 2.2.6，黎曼面的基本群 $\pi_1(X;x)=\Aut(\tilde X/X)^{\mathrm{op}}$ 对于 $X$ 的
复解析覆盖范畴所起的作用如同 Galois 群。克莱因曲面 $(X,\sigma)$ 的基本群（其定义见
本节定义 2.2.10）对于以 $(X,\sigma)$ 为底的实覆盖映射范畴起着同样的作用。下文中用
$\Aut^{\pm}(X)$ 记给定黎曼面 $X$ 的全纯与反全纯变换之群（它确实是群：若 $X$ 带有实
结构 $\sigma$，则它同构于半直积 $\Aut(X)\rtimes_\sigma\Z/2\Z$，其中
$\Z/2\Z=\{\Id_X,\sigma\}$ 通过 $\sigma\curvearrowright f=\sigma\circ f\circ\sigma^{-1}$
作用在 $X$ 的全纯变换群 $\Aut(X)$ 上，见习题 2.2.23）。$(X,\sigma)$ 的基本群将是
$(\Aut^{\pm}(\tilde X))^{\mathrm{op}}$ 的子群，但不总是半直积，而且无论如何都不典范地
是半直积（例如命题 2.2.13 中定义的半直积结构依赖于基点的选取\footnote{事实上它只依赖于
包含该实点的 $X^\sigma$ 的连通分支，见习题 2.2.31。}）。我们先研究 $(X,\sigma)$ 的
覆盖范畴，并为其建立 Galois 理论。

\begin{definition}[实覆盖及其同态]\label{def:2.2.7}
克莱因曲面 $(X,\sigma)$ 的\emph{实覆盖}是克莱因曲面同态
$q:(Y,\tau)\to(X,\sigma)$（即 $Y$ 与 $X$ 之间的全纯映射，满足 $q\circ\tau=\sigma\circ q$），
且它同时是拓扑覆盖映射。
\end{definition}

$X$ 的实覆盖之间的\emph{同态}是克莱因曲面同态 $\phi:(Y,\tau)\to(Y',\tau')$，且它同时
是 $X$ 的覆盖之间的同态（即满足 $q'\circ\phi=q$ 的全纯映射，其中
$q:(Y,\tau)\to(X,\sigma)$、$q':(Y',\tau')\to(X,\sigma)$）。映射
$\phi:(Y,\tau)\to(Y',\tau')$ 是 $(X,\sigma)$ 的实覆盖同态，等价于图
\[
\begin{tikzcd}[column sep=large, row sep=large]
Y \arrow[rr, "\phi"] \arrow[dd, "\tau"'] \arrow[dr, "q"] & & Y' \arrow[dd, "\tau'"] \arrow[dl, "q'"']\\
 & X \arrow[dr, "\sigma"'] & \\
 Y \arrow[rr, "\phi"'] \arrow[ur, "q"'] & & Y' \arrow[ul, "q'"']
\end{tikzcd}
\]
的可交换性，或更紧凑地写成
\[
\begin{tikzcd}[column sep=large, row sep=large]
(Y,\tau) \arrow[rr, "\phi"] \arrow[dr, "q"'] & & (Y',\tau') \arrow[dl, "q'"]\\
 & (X,\sigma) &
\end{tikzcd}
\]

一个自然的问题是：是否存在群 $\pi$，使得 $(X,\sigma)$ 的实覆盖范畴等价于离散左
$\pi$-集范畴，如定理 2.2.6 的模式那样。给定实覆盖 $q:(Y,\tau)\to(X,\sigma)$，由命题
2.2.5 我们已知 $\pi_1(X;x)$ 作用在 $q^{-1}(\{x\})$ 上：若 $y$ 是
$q^{-1}(\{x\})$ 中任一元素，则存在唯一的映射
$t_y:(\tilde X,\tilde x)\to(Y,y)$ 使得 $q\circ t_y=p$，其中
$p:(\tilde X,\tilde x)\to(X,x)$ 是 $X$ 的一个固定万有覆盖。现在，由于
$q:(Y,\tau)\to(X,\sigma)$ 是克莱因曲面同态，有离散集之间的双射
\[
\tau|_{q^{-1}(\{x\})}:\ q^{-1}(\{x\})\xrightarrow{\ \simeq\ }q^{-1}(\{\sigma(x)\}),
\]
而这个双射可用来定义比 $\pi_1(X;x)$ 严格更大的群在 $q^{-1}(\{x\})$ 上的作用。事实上，
假定 $f:\tilde X\to\tilde X$ 是 $\tilde X$ 的自微分同胚，它覆盖的不是 $\Id_X$，而是
实结构 $\sigma$，即
\[
\begin{tikzcd}[column sep=large, row sep=large]
\tilde X \arrow[r, "f"] \arrow[d, "p"'] & \tilde X \arrow[d, "p"]\\
X \arrow[r, "\sigma"'] & X
\end{tikzcd}
\]
此时
\[
q\bigl((t_y\circ f)(\tilde x)\bigr)
=\bigl((q\circ t_y)\circ f\bigr)(\tilde x)
=(p\circ f)(\tilde x)=(\sigma\circ p)(\tilde x)=\sigma(x),
\]
故 $(t_y\circ f)(\tilde x)\in q^{-1}(\{\sigma(x)\})$，从而
\begin{equation}\label{eq:2.13}
f\curvearrowright y:=(\tau^{-1}\circ t_y\circ f)(\tilde x)\in q^{-1}(\{x\}).
\end{equation}
换言之，$Y$ 上与到 $(X,\sigma)$ 的投影相容的实结构 $\tau$ 的存在性，蕴涵着下述群
左作用于 $q$ 在 $x$ 处的纤维：
\begin{equation}\label{eq:2.14}
\pi^R_1((X,\sigma);x):=
\bigl\{f:\tilde X\to\tilde X\ \bigm|\ \exists\sigma_f\in\{\Id_X;\sigma\},\
p\circ f=f\circ\sigma_f\bigr\}^{\mathrm{op}}
\subset\Aut^{\pm}(\tilde X)^{\mathrm{op}}
\end{equation}
\begin{annote}
式 (2.14) 中的条件 $p\circ f=f\circ\sigma_f$ 照底本原样排印。按各映射的定义域与值域，
此处应为 $p\circ f=\sigma_f\circ p$（同页的 (2.15) 与 §2.3 的 (2.20) 即如此写）；译文不改动
底本，谨此注明。
\end{annote}

下述引理（命题 2.2.8）将证明：这个群确实以
$\pi_1(X;x)$ 为其满足 $\sigma_f=\Id_X$ 的元素 $f$ 所成子群（特别地，这是正规子群：
因为若 $g$ 覆盖 $\mu\in\{\Id_X,\sigma\}$，则 $g^{-1}\circ f\circ g$ 覆盖
$\mu^{-1}\circ\Id_X\circ\mu=\Id_X$）。$\pi^R_1((X,\sigma);x)$ 的另一种定义是
\begin{equation}\label{eq:2.15}
\pi^R_1((X,\sigma);x)=\bigcup_{\mu\in\{\Id_X;\sigma\}}
\bigl\{f:\tilde X\to\tilde X\ \bigm|\ p\circ f=\mu\circ p\bigr\}.
\end{equation}
需要验证的有几件事，我们把它们分成两个结果来陈述。
